
Soit $M$ une catégorie de motifs géométriques, [de type fini,] sur le
préschéma $S$ de type fini sur $\mathbb{Z}$. Pour simplifier, supposons que
les degrés de $M$ soient pairs (sinon, il sera possible, pour
[\uncertain{néanmoins}] \struck{\ill{}} [traiter] la question de la
distribution des Frobenius, de remplacer $M$ par $M^{(\mathrm{pair})}$).
Tout \uncertain{choix} d'un \ill{} permet de définir des pts fermés
$s \in S$ \ill{} « $\mathbb{R}$-points motiviques » \struck{de $M$}
\ill{} \ill{}. Pour fixer les idées, \ill{} \ill{} un pt
$\mathbb{R}$-\uncertain{valué} \ill{} \ill{} \uncertain{algébrique}
$G_{\mathbb{R}}$ sur $\mathbb{R}$ de $M$, \ill{} \ill{} \ill{}.

Il y a lieu de \struck{\ill{}} [\uncertain{démultiplier}] les points
[fermés] de $S$, en considérant les couples $s' = (s, \xi)$, où
$s \in S^{0}$, et \ill{} \struck{\ill{}} \uncertain{géométrique} \ill{}
est du type
\[
  i : \mathbb{G}_{m,\mathbb{C}} \to G(s)_{\mathbb{C}}
\]
\ill{} Galois motivique associé à $M_s$, --- $\mathbb{C}(s)$ désigne la
\ill{} \ill{} \ill{} \uncertain{complexe}. À tout $s'$ \ill{}

\marginal{l'existence d'un \ill{} \uncertain{désigne} \ill{}
\uncertain{généraux}}
\note{note oblique dans la marge gauche, le dernier mot souligné}

\struck{Alors} Soit $S'$ l'ensemble des $\mathbb{R}$-points motiviques de
$M$, \struck{et} \ill{} \ill{} déterminé à isom.\ près \struck{\ill{}}
\[
  \xi(s') \in H^{1}(\mathbb{R}, G^{0u}_{\mathbb{R}})
\]
où $G^{0u}$ est le noyau de $G^{0}_{\mathbb{R}}
\xrightarrow{\ \varepsilon\ } \mathbb{G}_m$. De plus, si
$G_{\mathbb{R}}^{\xi(s')}$ est le « \uncertain{tordu} » \ill{} de
$G_{\mathbb{R}}$ par torsion par $\xi(s')$, on obtient \ill{} un
homomorphisme \ill{}
\[
  u(s') : G(s')_{\mathbb{R}} = G(s)_{\mathbb{R}} \longrightarrow
  G_{\mathbb{R}}^{\xi(s')}
\]
compatible avec $\varepsilon_{\mathbb{R}}$, $i_{\mathbb{R}}$, et déterminé
en fait mod \struck{automorphismes intérieurs} [\ill{}] par un élément de
$G^{0u}_{\mathbb{R}}(\mathbb{R})$.
\note{l'exposant de $G^{0u}$ se lit mal : un rond suivi d'une boucle, lu
« $0u$ » ; on pourrait aussi y voir « $0n$ ». Un premier $G$ en tête de la
formule $\xi(s')$ est surchargé}

